%%%----------------------------------------------------------------------------------------
%	SECTION TITLE
%----------------------------------------------------------------------------------------
\vspace{-3mm}
\cvsection{Experience}

%----------------------------------------------------------------------------------------
%	SECTION CONTENT
%----------------------------------------------------------------------------------------

\begin{cventries}

%------------------------------------------------

\cventry
{Staff Engineer II (Engineering Physicist II)}
{SLAC National Laboratory}
{Palo Alto, CA}
{April 2022 - Present}
{
\begin{cvitems}
\item {Led the recommissioning of the bunch length diagnostic system after six years offline, restoring a critical beam diagnostic and enabling reliable, high-resolution measurements for accelerator operations and research.}
\item {Developed MATLAB-based tools to process and visualize beam diagnostics data across accelerator sectors, including magnet mover controls and slice emittance and energy spread analysis, improving data accessibility and analytical capability for operations and research teams.}
\item {Supported accelerator operations and troubleshooting, contributing to stable beam delivery and system optimization during FACET-II runs.}
\item {Led the restoration of the FACET positron beamline (inactive since 2016), returning RF, magnet, vacuum, and water cooling systems to full operational status.}
\item {Directed cross-disciplinary engineering efforts across vacuum, electrical, and cooling systems to re-establish beamline functionality and ensure safe, reliable operation.}
\item {Engineered electrical configurations to enforce robust Lockout-Tagout (LOTO) conditions, significantly improving safety and maintenance accessibility.}
\end{cvitems}
}
\vspace{1.5mm}
%------------------------------------------------    
\cventry
{Engineering Physicist I, US-HL-AUP Superconducting R\&D Engineer \& Electrical Quality Assurance}
{Fermilab}
{Batavia, IL}
{April 2015 - March 2022}
{
\begin{cvitems}
\item {Managed daily fabrication workflow for superconducting magnet coils, supporting pre-production (24 units) and full production (50 units), optimizing production rates and addressing process inefficiencies in collaboration with technicians to meet delivery targets.}
\item {Led electrical quality control for Nb$_3$Sn superconducting magnet systems in the US HiLumi-LHC AUP, improving fabrication reliability and reducing defect rates during early-stage production.}
\item {Led the successful design and fabrication of R\&D test coils used to investigate failure modes in epoxy-injected superconducting cables.}
\item {Led electromechanical failure analyses of prototype and production Nb$_3$Sn magnet coils to diagnose quench causes, informing time-critical production decisions.}
\item {Initiated and led a review to enhance and update electrical design parameters applied to AUP magnet coils during early fabrication stages.}
\item {Designed and implemented a nitrogen (N$_2$) distribution system for furnace operations to control oxygen during Nb$_3$Sn coil processing, improving process reliability and reducing technician intervention.}
\item {Developed a MATLAB GUI to analyze thermocouple temperature profiles against heat treatment recipes, identifying process control discrepancies and supporting validation of furnace performance.}
\item {Presented heat treatment and electrical quality control (QC) data on a per-coil basis to support weekly readiness reviews, informing decisions on progression to subsequent fabrication stages.}
\item {Supervised electromechanical technicians on specialized R\&D projects for cold-magnet coil fabrication efforts.}
\item {Worked closely with physicists to optimize and improve the heat treatment process for AUP coils to achieve desired superconducting performance goals based on laboratory testing and analysis.}
\item {Executed final electrical qualification of AUP superconducting coils prior to shipment to CERN, verifying compliance with specifications and ensuring readiness for integration.}
\end{cvitems}
}
\vspace{3mm}
%------------------------------------------------

\cventry
{Accelerator Systems \& Operations}
{Fermilab}
{Batavia, IL}
{Feb 2013 - April 2014}
{
\begin{cvitems}
\item {Coordinated accelerator operations to initiate, optimize, and maintain beamline performance in support of laboratory and university-based research programs.}
\item {Operated, tuned, and maintained the H$^-$ plasma source to ensure stable beam delivery across the accelerator complex.}
\item {Contributed to the design and implementation of an air diversion system for the NuMI horn power stripline, improving thermal management and operational reliability for high-power beam delivery.}
\end{cvitems}
}
\vspace{6mm}
%------------------------------------------------
\cventry
{Graduate Research Assistant}
{Georgia Institute of Technology}
{Atlanta, GA}
{Dec 2013 - April 2014}
{
\begin{cvitems}
\item {Participated in an in-depth study to understand plasma boundary interactions and wall effects.}
\item {Studied plume divergence behavior of the P5 5 kW Hall thruster.}
\item {Assisted in the design of a glass injection tube for a helicon ion thruster.}
\item {Aided in the fabrication of several hollow cathodes and plasma probes for testing.}
\end{cvitems}
}
\vspace{6mm}
%------------------------------------------------
\cventry
{Active Aeroelasticity and Structures Lab - Undergraduate Research Assistant}
{University of Michigan}
{Ann Arbor, MI}
{Dec 2012 - Dec 2013}
{
\begin{cvitems}
\item {Conducted aeroelasticity research on highly flexible aircraft, supporting design and testing of UAV systems.}
\item {Designed and validated a ventral fin using CFD and wind tunnel testing (2'×2' and 5'×7') to improve yaw stability and control authority.}
\item {Developed a sensor-based system using gyroscopes to model and measure aircraft structural deformation in flight.}
\item {Fabricated and maintained structural components for the X-HALE UAV, supporting iterative design improvements.}
\item {Contributed to full-scale assembly of an Aerial Test Vehicle (ATV), including composite wing layup and manufacturing.}
\end{cvitems}
}
%------------------------------------------------
%------------------------------------------------
\end{cventries}